\section{Discussion}
\label{sec:discussion}

\subsection{What has been established}
\label{sec:established}

Three results stand on the evidence presented. First, the framework computes
the surface gravity field correctly: given identical inputs it reproduces an
independent Werner--Scheeres calculation facet by facet, with an RMS difference
of half a degree and $99.96\%$ of facets agreeing to within five degrees
(Section~\ref{sec:ryugu}), and given only a shape model and an assumed density
it recovers the published global mean, latitudinal structure, geopotential
pattern and density dependence for Bennu (Sections~\ref{sec:bennu}
and~\ref{sec:densitysweep}).

Second, the area-weighted mean slope retains a substantial mesh dependence over
the tested refinement range. On the three bodies with the most detailed shape
models it had not converged at the finest mesh available, still moving by
$2.6$--$8.2\%$ over the last refinement step (Section~\ref{sec:slopesens}).
This is a property of the diagnostic, and the smoothing experiment identifies
its cause as relief that coarse meshes fail to resolve.

Third, the normalized geopotential dispersion computed from the same field on
the same meshes shows clear practical convergence, changing by at most $0.6\%$
over the same final step
and by nothing at all on two bodies (Section~\ref{sec:varrobust}). The
distinction follows from what the two quantities depend on: a volume integral
of the mass distribution in one case, the orientation of an individual facet in
the other.

None of this makes the slope a less useful quantity---it remains the direct
measure of whether material will move---but it does mean that a slope value is
a statement about a particular mesh, and must be reported as one.

\subsection{The assumed density dominates the error budget}
\label{sec:densitydominates}

It would be a mistake to leave the impression that mesh resolution is the
principal difficulty facing shape-based surface gravity. For any body whose
mass has not been measured, the assumed bulk density is a far larger source of
uncertainty.

The Bennu sweep of Table~\ref{tab:sweep} makes the comparison concrete. Varying
the assumed density across the range that would be plausible for a carbonaceous
rubble pile in the absence of a mass determination, from $0.85$ to
$1.194\ \mathrm{g\,cm^{-3}}$, moves the mean slope by $10.1^{\circ}$. Mesh
resolution, on the bodies where we measured it, moves it by $2$--$3.5^{\circ}$.
The density assumption is thus the dominant term by a factor of three or more,
and it acts as a systematic shift rather than as scatter.

This ordering has a practical consequence. For a body known only from the
ground, refining the shape model buys less than constraining the mass would,
and the two are not substitutes. It also reinforces the argument of
Section~\ref{sec:xvimplications}: in the one controlled case available, a
coarse ground-based model reproduced the mean slope to within one per cent and
the geopotential dispersion to within eight per cent, while the density
assumption contributes considerably more, so the shape model is rarely the
limiting factor for global quantities.

\subsection{The uniform-density assumption}
\label{sec:uniformdensity}

Every field in this paper assumes a homogeneous interior. That assumption is
exact for none of the bodies considered, and it is worth being explicit about
where it bites and where it does not.

It does not affect the resolution results of Sections~\ref{sec:slopesens}
and~\ref{sec:varrobust}. Those hold density fixed and vary only the mesh, so
whatever error the homogeneity assumption introduces is common to every point
of a sequence and cancels from $S$ and $C$.

It does affect the absolute fields, and Itokawa is the clearest case. Its head
and body are inferred to differ in density by roughly a quarter, with the head
the denser of the two \citep{kanamaru2019}, a conclusion supported
independently by the YORP spin-up analysis of \citet{lowry2014}. A homogeneous
calculation must therefore misplace the local gravity direction near the neck,
where the two lobes contribute comparably. The comparison with the Ryugu
archive in Section~\ref{sec:ryugu} does not test this, since that reference is
itself computed for a uniform interior; it validates the solver, not the
physical assumption.

For bodies whose interiors are close to homogeneous---which the near-coincidence
of the centre of figure and centre of mass suggests for Eros---the assumption is
mild. For contact binaries and bilobed objects it is not, and slopes computed
near the junction should be treated with corresponding caution.

\subsection{Toward interior structure}
\label{sec:interior}

The limitation just described is also an opportunity, and it motivates the
companion work.

If a surface has been shaped over long periods by downslope transport, it should
tend toward an equipotential of the true gravity field, including whatever
density inhomogeneity that field contains. A trial interior density
distribution that is closer to the truth should therefore produce a smaller
spread of geopotential over the surface. This is the reasoning by which
\citet{kanamaru2019} inferred a denser head for Itokawa, and it converts the
geopotential dispersion from a descriptive statistic into an estimator.

Such an inversion places a specific demand on the diagnostic used: its variation
with the assumed density distribution must not be swamped by its variation with
the mesh. Section~\ref{sec:varrobust} establishes that $\potvarn$ satisfies
this, and Section~\ref{sec:slopesens} that the mean slope does not. That is the
main reason the distinction developed in this paper matters beyond questions of
reporting practice.

Two further ingredients are already in place. The linearity of the field in
density (Section~\ref{sec:solver}) means that a body partitioned into regions
of distinct density can be evaluated by superposing unit-density solutions, so
that scanning a density model costs one solver call per region rather than one
per trial. And Paper~I's bound $GM/\RA^{2}\le\gmean$ \citep{paper1} holds for
any interior mass distribution, providing a constraint that survives the
homogeneity assumption. Developing these into a working estimator, validating
it against the independently constrained cases, and establishing what can and
cannot be recovered from surface data alone is reserved for future work and is
not attempted here.

\subsection{On verification practice}
\label{sec:verifpractice}

One methodological point deserves emphasis beyond its role in this work.

The natural first test of a surface gravity code is a homogeneous sphere, which
must return zero slope everywhere. That test is necessary but weak: a
non-rotating sphere returns zero slope for either sign of the centrifugal term,
and so cannot detect an error in the rotating frame---which is precisely the
part of the calculation that distinguishes a small-body code from a
straightforward Newtonian one. The uniformly rotating sphere of
Eq.~\eqref{eq:rotsphere} does discriminate, and sharply: at
$\omegarot^{2}R/g = 0.5$ the two signs differ by more than seven degrees at
mid-latitude.

The cost of the stronger test is negligible---it is a closed-form expression and
a few lines of code, both provided with this work---and we suggest it be adopted
routinely. The same applies to checking the potential against the closed-form
ellipsoid solution of Eq.~\eqref{eq:ellipsoid}, which exercises the
gravitational term with the rotation switched off. Between them the two tests
cover every term in Eqs.~\eqref{eq:Ug}--\eqref{eq:geopot}, and neither depends
on any published data set.

\subsection{Limitations}
\label{sec:limitations}

The scope of what has been demonstrated should be stated plainly.

The per-facet validation rests on a single body. Ryugu is the only object for
which per-facet gravitational products computed on a publicly available shape
model at a stated density and spin are distributed, and it is therefore the only
case in which the implementation could be checked below the level of aggregate
statistics. That test is stringent, but it is one test.

The cross-version comparison rests on a single body as well. Of the two
asteroids in our set with independent shape models, only Itokawa passes the
prerequisite that the models describe a body of the same size
(Section~\ref{sec:kleopatraxv}). Whether the division between recoverable bulk
descriptors and unrecoverable distribution tails holds generally is not
established by one case.

The resolution study covers five bodies and a range of roughly $2\,000$ to
$42\,000$ facets. Whether the geopotential dispersion continues to hold its plateau
at $10^{5}$ or $10^{6}$ facets is untested, and the highest-resolution models
now available for Bennu and Ryugu would allow that range to be extended.

Finally, the slope of Eq.~\eqref{eq:slope} is one of several definitions in use.
It is the definition used by the Ryugu reference products, which is what makes
the comparison of Section~\ref{sec:ryugu} meaningful. Studies that define slope
relative to a geopotential-based reference surface rather than to the local
facet normal will obtain different values, particularly in equatorial regions of
rapidly rotating bodies, and the numbers here should not be compared with those
without accounting for the difference.
